\section{Representing semantic rules} \label{sec:semantics_lean}

The basic semantics of ShEx formalizes the \lean{Satisfies} ($\vdash$) and \lean{Matches} ($\vDash$) 
 relations of Fig. 2 in~\cite{BGP17}. 
both parameterized by a typing t. 
The two relations are mutually recursive, 
 a shape expression's neighbourhood description invokes \lean{Matches} 
 on the node's outgoing triples, 
 and a triple expression's shape reference invokes \lean{Satisfies} 
 indirectly, through membership in \lean{t} 
 so they are declared together as a single mutual block, 
 \lean{Satisfies} and \lean{Matches}, 
 mirroring the paper's joint Fig. 2a/2b presentation.

Although \lean{Satisfies} and \lean{Matches} look like a classic mutually 
 inductive derivability relation, 
  they are encoded as ordinary Prop-valued recursive functions 
  rather than a mutual inductive family. 
  This is necessary because the premise in the rule involving negation 
  is a negative occurrence of \lean{Satisfies}, 
  which Lean's positivity checker rejects for an inductive definition 
  but accepts for a function.
  
\lean{Satisfies} and \lean{Matches} are defined case by case. Each equation corresponds to one (or a disjunction of two) rules.

\begin{leanlisting}
mutual
def Satisfies (G : Graph) (t : Typing) (n : RDFNode) : ShExpr → Prop
  | .valueDescr V => InV n V ∨ BNodeInV n V
  | .neigDescr E => Matches G t (neigh G n) E
  | .and S1 S2 => Satisfies G t n S1 ∧ Satisfies G t n S2
  | .or S1 S2 => Satisfies G t n S1 ∨ Satisfies G t n S2
  | .not S => ¬ Satisfies G t n S

def Matches (G : Graph) (t : Typing) (N : Finset Triple) : TExpr → Prop
  | .tPattern P L =>
      ∃ subj pred obj, N = {(⟨subj, pred, obj⟩ : Triple)} ∧ pred ∈ P ∧ t.has obj L
  | .empty => N = ∅
  | .someOf E1 E2 => Matches G t N E1 ∨ Matches G t N E2
  | .eachOf E1 E2 =>
      ∃ N1 N2, DisjointUnion N N1 N2 ∧
               Matches G t N1 E1 ∧
               Matches G t N2 E2
  | .rep E min max =>
      ∃ Ns : List (Finset Triple),
        DisjointUnionList N Ns ∧
        min ≤ Ns.length ∧
        (∀ m', max = some m' → Ns.length ≤ m') ∧
        ∀ Ni ∈ Ns, Matches G t Ni E
end
\end{leanlisting}

\lean{InV} is a proposition that passes if \lean{n} 
 is an RDF Node in \lean{V}, 
 \lean{BNodeInV} passes if \lean{n} 
 is a blank node in \lean{V}, 
 and \lean{neigh} is a function that returns the set of triples in \lean{G}
 that have \lean{n} as subject. \lean{DisjointUnion} and \lean{DisjointUnionList} are defined as follows:

\begin{leanlisting}
def DisjointUnion (N N1 N2 : Finset Triple) : Prop :=
  Disjoint N1 N2 ∧ N = N1 ∪ N2

def DisjointUnionList (N : Finset Triple) 
                      (Ns : List (Finset Triple)) : Prop :=
  Ns.Pairwise Disjoint ∧ N = Ns.foldr (· ∪ ·) ∅
\end{leanlisting}

  

